% Generated deterministically from the audited Markdown source.
% Responsible author: Carptopus.
\documentclass[11pt]{article}
\usepackage[a4paper,margin=27mm]{geometry}
\usepackage[T1]{fontenc}
\usepackage[utf8]{inputenc}
\usepackage{lmodern}
\usepackage{microtype}
\usepackage{amsmath,amssymb,amsthm,mathtools}
\usepackage{needspace}
\usepackage{xcolor}
\usepackage[colorlinks=true,linkcolor=blue!55!black,citecolor=blue!55!black,urlcolor=blue!65!black]{hyperref}
\usepackage{xurl}
\usepackage{fancyhdr}

\newtheorem{theorem}{Theorem}[section]
\newtheorem{lemma}[theorem]{Lemma}

\pagestyle{fancy}
\fancyhf{}
\fancyhead[L]{Carptopus}
\fancyhead[R]{Affine stabilizers of o-monomial value sets}
\fancyfoot[C]{\thepage}
\setlength{\headheight}{14pt}
\setlength{\parindent}{0pt}
\setlength{\parskip}{0.44em}
\setlength{\emergencystretch}{4em}
\urlstyle{same}

\title{Affine stabilizers of two o-monomial value sets}
\author{Carptopus\\\href{mailto:carptopus@163.com}{carptopus@163.com}}
\date{4 September 2026}

\hypersetup{
  pdftitle={Affine stabilizers of two o-monomial value sets},
  pdfauthor={Carptopus},
  pdfsubject={Trivial affine stabilizers for the overline-Segre and Glynn I value sets},
  pdfkeywords={o-monomial, o-polynomial, hyperoval, finite field, affine stabilizer, value-set polynomial}
}

\begin{document}
\maketitle

\begin{center}
\fcolorbox{black!35}{black!4}{\parbox{0.88\textwidth}{\centering\small
Public Beta manuscript, version 0.1-beta. The mathematical proof and complete
manuscript have passed project-internal zero-trust audits. External mathematical review
and formal peer review are pending.}}
\end{center}

\begin{abstract}


Let \(q=2^m\), where \(m\ge 5\) is odd, and let

\[
J_e=\{x^e+x:x\in\mathbb F_q\}.
\]

Ding and Tang asked whether the stabilizer of \(J_e\) in the one-dimensional
general affine group is trivial for every nontranslation o-monomial \(x^e\).
Their argument settled the known families whose exponents are sums of two powers
of two and left two known families untreated. We prove triviality for both of
those families: the \(\overline{\mathrm{Segre}}\) family \(e=q-6\) and the
Glynn I family \(e=3\cdot2^{(m+1)/2}+4\). The first proof uses a general
criterion involving the first three coefficients of the value-set polynomial.
The second locates the first nonzero coefficient by expressing a circulant
determinant as a parity count of directed cycle covers. Consequently, the
conjectured stabilizer property holds for every currently known nontranslation
o-monomial family. This does not classify hypothetical unknown o-monomials.

\end{abstract}

\section{Introduction}


An o-polynomial over \(\mathbb F_q\), with \(q\) even, represents a hyperoval in
the Desarguesian projective plane. When the o-polynomial is a monomial \(x^e\),
the map \(x\mapsto x^e+x\) is two-to-one, and hence its image \(J_e\) has size
\(q/2\).

Ding and Tang \cite{ding-tang} use the orbit of \(J_e\) under the general affine group

\[
\mathrm{GA}_1(\mathbb F_q)
=\{x\mapsto ux+v:u\in\mathbb F_q^*,\ v\in\mathbb F_q\}
\]

to construct \(3\)-designs. The orbit size, and therefore the design parameters,
depend on the affine stabilizer of \(J_e\). Their Conjecture 2 predicts that this
stabilizer is trivial for every nontranslation o-monomial. They proved the claim
for the known two-power-sum exponents and explicitly left the
\(\overline{\mathrm{Segre}}\) and Glynn I families open \cite[after Corollary~30]{ding-tang}.

This note treats both remaining known families. Its two main ingredients are:

\begin{enumerate}
\item a low-coefficient criterion that detects a trivial affine stabilizer for a
half-size subset of \(\mathbb F_q\); and
\item a cycle-cover calculation that determines the first nonzero coefficient of
the Glynn I value-set polynomial.

\end{enumerate}

Kölsch and Kyureghyan \cite{kolsch-kyureghyan} subsequently established an equivalence between the
classification of o-monomials and that of two-to-one binomials. Their result
clarifies the classification boundary but does not determine the affine
stabilizer of the fixed image set \(J_e\). Thus the conclusion here concerns all
currently known nontranslation families, not all possible future o-monomials.

All polynomial coefficients below lie in \(\mathbb F_q\). Since the
characteristic is two, signs in elementary symmetric functions and determinants
disappear. Throughout the remainder of the paper, \(q=2^m\) with odd
\(m\ge5\), unless explicitly stated otherwise.

Maschietti \cite{maschietti} relates the nonzero part of these value sets to cyclic difference
sets. Our proofs require finer value-set coefficient information rather than
only the difference-set parameters.

\section{A value-set coefficient criterion}


Let \(J\subseteq\mathbb F_q\) have size \(k=q/2\), and write its monic root
polynomial as

\[
P_J(T)=\prod_{y\in J}(T+y)
=T^k+c_1T^{k-1}+c_2T^{k-2}+c_3T^{k-3}+\cdots.
\]

\Needspace{0.18\textheight}
\begin{lemma}
\label{lem:criterion}

Suppose that \(m\ge5\) is odd and

\[
c_1=c_2=0,\qquad c_3\ne0.
\]

Then the stabilizer of \(J\) in \(\mathrm{GA}_1(\mathbb F_q)\) is trivial.

\end{lemma}

\begin{proof}

Assume that \(uJ+v=J\), where \(u\ne0\). Equality of the corresponding monic
root polynomials gives

\[
P_J(T)=u^kP_J\!\left(\frac{T+v}{u}\right).
\tag{2.1}
\]

Because \(k\) is a power of two, \((T+v)^k=T^k+v^k\). Comparing the
coefficient of \(T^{k-3}\) in (2.1) yields \(c_3=u^3c_3\), so \(u^3=1\).
Since \(m\) is odd, \(\gcd(3,q-1)=1\), and therefore \(u=1\).

Now \(P_J(T)=P_J(T+v)\). Comparing the coefficient of \(T^{k-4}\), the two
copies of \(c_4\) cancel and the term \(c_3(T+v)^{k-3}\) contributes
\(c_3v\), because \(k-3\) is odd. Hence \(v=0\). The assumption \(m\ge5\)
ensures that this comparison is above the constant term.
\end{proof}


\section{The overline-Segre family}


Set \(e=q-6\) and

\[
J=\{x^{q-6}+x:x\in\mathbb F_q\}.
\]

The monomial \(x^{q-6}\) is the
\(\overline{\mathrm{Segre}}\) o-monomial in the parameter range under
consideration \cite[Theorem~12 and the following identities]{ding-tang}. Thus
\(x\mapsto x^{q-6}+x\) is two-to-one and \(|J|=q/2\).

\Needspace{0.18\textheight}
\begin{theorem}
\label{thm:bar-segre}

For every odd \(m\ge5\), the stabilizer of \(J\) in
\(\mathrm{GA}_1(\mathbb F_q)\) is trivial.

\end{theorem}

\begin{proof}

Write \(k=q/2\) and

\[
P(T)=\prod_{y\in J}(T+y)
=T^k+c_1T^{k-1}+c_2T^{k-2}+c_3T^{k-3}+\cdots.
\]

The two-to-one property gives the multiset identity

\[
P(T)^2=\prod_{x\in\mathbb F_q}(T+x^{q-6}+x).
\tag{3.1}
\]

For \(x\ne0\), put \(z_x=x+x^{-5}\), and let \(E_r\) be the \(r\)-th
elementary symmetric function of the \(z_x\), indexed by
\(\mathbb F_q^*\). Equation (3.1) implies

\[
c_1^2=E_2,\qquad c_2^2=E_4,\qquad c_3^2=E_6.
\tag{3.2}
\]

Expand a term of \(E_r\) by choosing \(x\) from \(a\) factors and \(x^{-5}\)
from the remaining \(r-a\) factors. Multiplication of every selected field
element by \(t\in\mathbb F_q^*\) multiplies the resulting subsum by

\[
t^{\,6a-5r}.
\]

As this multiplication permutes the indexing subsets, that subsum is zero
unless \(q-1\mid 6a-5r\). For \(r=2\) and \(r=4\), all possible weights are
nonzero and have absolute value less than \(q-1\ge31\), so \(E_2=E_4=0\).
For \(r=6\), only \(a=5\) has weight zero. Consequently,

\[
E_6=\sum_{b\in\mathbb F_q^*}b^{-5}
e_5(\mathbb F_q^*\setminus\{b\}).
\]

Since, for a formal indeterminate \(Z\),

\[
\prod_{x\in\mathbb F_q^*}(1+Zx)=1+Z^{q-1},
\]

division by \(1+Zb\) shows that
\(e_5(\mathbb F_q^*\setminus\{b\})=b^5\). Hence

\[
E_6=\sum_{b\in\mathbb F_q^*}1=q-1=1
\]

in characteristic two. By (3.2), \(c_1=c_2=0\) and \(c_3=1\). Lemma~\ref{lem:criterion}
completes the proof.
\end{proof}


\section{The Glynn I family}


Let

\[
e=3\cdot2^{(m+1)/2}+4
\]

be the Glynn I exponent. Put

\[
h=\frac{m-1}{2},\qquad A=2^h,\qquad
N=q-1=2A^2-1,\qquad D=e-1=3(2A+1),
\]

and let \(k=q/2=A^2\). Define

\[
J_e=\{x^e+x:x\in\mathbb F_q\},\qquad
P(T)=\prod_{y\in J_e}(T+y)
=T^k+c_1T^{k-1}+c_2T^{k-2}+\cdots.
\]

Again the o-monomial property gives

\[
P(T)^2=\prod_{x\in\mathbb F_q}(T+x^e+x).
\tag{4.1}
\]

The o-monomial property used here is Glynn's construction \cite{glynn}, recorded in
\cite[Theorem~13]{ding-tang}.

\Needspace{0.18\textheight}
\begin{lemma}
\label{lem:glynn-coefficient}

The first nonzero coefficient after the leading term of \(P\) is \(c_j=1\),
where

\[
j=
\begin{cases}
2^{(m-3)/2},&m\equiv1\pmod4,\\[4pt]
3\cdot2^{(m-3)/2}-1,&m\equiv3\pmod4.
\end{cases}
\tag{4.2}
\]

Moreover, \(\gcd(j,N)=1\). In the case \(m\equiv1\pmod4\), one also has

\[
c_{j+1}=\cdots=c_{2j-1}=0.
\tag{4.3}
\]

\end{lemma}

\begin{proof}

Let \(E_R\) denote the \(R\)-th elementary symmetric function on the multiset
\((x^e+x)_{x\in\mathbb F_q^*}\). From (4.1),

\[
c_i^2=E_{2i}.
\tag{4.4}
\]

In an expanded term of \(E_R\), suppose \(b\) factors contribute \(x^e\) and
the other \(R-b\) factors contribute \(x\). Scaling the selected field elements
shows that the corresponding subsum vanishes unless

\[
N\mid R+bD,\qquad 0\le b\le R.
\tag{4.5}
\]

The least positive \(R\) satisfying (4.5) is

\[
(R_0,b_0)=
\begin{cases}
\left(A,(A-1)/3\right),&h\text{ even},\\[4pt]
\left(3A-2,(7A-5)/3\right),&h\text{ odd}.
\end{cases}
\tag{4.6}
\]

For completeness, if \(h\) is even, then \(A\equiv1\pmod3\),
\(N\equiv A\pmod D\), and the bound \(R(1+D)<3N\) for \(R<A\) leaves only
the residues \(A\) and \(2A\). If \(h\) is odd, then \(A\equiv2\pmod3\),
\(N\equiv3A+1\pmod D\), and \(R(1+D)<9N\) for \(R<3A-2\). The least positive
residues of \(tN\pmod D\), \(1\le t\le8\), first enter the permitted range at
\(t=7\), where the residue is \(3A-2\). The displayed \(b_0\) then follows by
solving \(R_0+b_0D=tN\). Also \(R_0<N\) for \(A\ge4\), and
\(\gcd(D,N)=1\), so this \(b_0\) is unique in \(0\le b\le R_0\).

It remains to show that the surviving subsum at \(R_0\) equals one. Let \(S\)
be the \(N\times N\) cyclic shift matrix. Its eigenvalues are the elements of
\(\mathbb F_q^*\), and therefore

\[
\prod_{x\in\mathbb F_q^*}(1+Ux^e+Vx)
=\det(I+US^e+VS).
\tag{4.7}
\]

The coefficient of \(U^bV^{R-b}\) is the parity of cycle covers having \(b\)
steps of size \(e\), \(R-b\) steps of size \(1\), and all remaining vertices
fixed. Minimality of \(R_0\) implies that a cover with \(R_0\) nonfixed edges
contains exactly one nontrivial cycle. No proper consecutive subword of a step
word can sum to zero modulo \(N\), again by minimality. Hence the number of such
cycles is

\[
\frac{N}{R_0}\binom{R_0}{b_0}.
\tag{4.8}
\]

One has \(\gcd(N,R_0)=1\): this is immediate when \(R_0=A\), while for
\(R_0=3A-2\) it follows from \(9N\equiv-1\pmod{R_0}\). Kummer's theorem gives

\[
v_2\!\binom{R_0}{b_0}=v_2(R_0).
\tag{4.9}
\]

For \(h\) even, \(R_0=A\) and \(b_0=(A-1)/3\) is odd, giving valuation \(h\).
For \(h\) odd, put \(c_0=R_0-b_0=(2A-1)/3\). The alternating binary
expansions give

\[
s_2(c_0)=\frac{h+1}{2},\qquad
s_2(b_0)=\frac{h+1}{2},\qquad
s_2(R_0)=h.
\]

Kummer's digit-sum formula therefore gives
\(v_2\binom{R_0}{b_0}=1=v_2(R_0)\). Thus (4.8) is odd, so \(E_{R_0}=1\).
Equations (4.4) and (4.6) yield (4.2), and the formulas also give
\(\gcd(j,N)=1\).

Finally suppose \(m\equiv1\pmod4\). Then \(R_0=A\) and \(j=A/2\). For
\(j<i<2j\), the degree \(R=2i\) lies strictly between \(A\) and \(2A\).
Using \(R(1+D)<7N\), the possible residues \(tN\pmod D\), \(1\le t\le6\),
are \(A,2A,\ldots,6A\), so none lies in this open interval. This proves (4.3).

\end{proof}


\Needspace{0.18\textheight}
\begin{theorem}
\label{thm:glynn}

For every odd \(m\ge5\), the stabilizer of the Glynn I value set \(J_e\) in
\(\mathrm{GA}_1(\mathbb F_q)\) is trivial.

\end{theorem}

\begin{proof}

Suppose \(uJ_e+v=J_e\). As in (2.1),

\[
P(T)=u^kP\!\left(\frac{T+v}{u}\right)
=\sum_{i\ge0}c_i u^i(T+v)^{k-i},\qquad c_0=1.
\tag{4.10}
\]

At degree \(T^{k-j}\), Lemma~\ref{lem:glynn-coefficient} gives \(1=u^j\). Since
\(\gcd(j,q-1)=1\), we have \(u=1\).

If \(m\equiv3\pmod4\), then \(j\) is odd. Comparing the coefficient of
\(T^{k-j-1}\) in \(P(T)=P(T+v)\) gives \(v=0\).

If \(m\equiv1\pmod4\), then \(j=A/2\). By (4.3), comparison at
\(T^{k-2j}\) leaves only

\[
\binom{k-j}{j}v^j=0.
\]

The lowest nonzero binary digit of \(k-j=A^2-A/2\) is precisely \(j\), so
Lucas' theorem gives \(\binom{k-j}{j}\equiv1\pmod2\). Hence \(v=0\) in this
case as well.
\end{proof}


\section{Consequence and scope}


Combining Theorems 3.1 and 4.2 with Ding and Tang's Lemma 29 proves the
conjectured trivial-stabilizer property for every nontranslation o-monomial
family currently listed in the classification literature. In particular, the
two known families explicitly left open in \cite{ding-tang} are now covered
\cite[Lemma~29 and after Corollary~30]{ding-tang}; \cite[Section~2]{kolsch-kyureghyan}.

This statement has a deliberate boundary. The classification of o-monomials is
itself open \cite{kolsch-kyureghyan}. We neither classify unknown families nor prove a theorem that
automatically covers every hypothetical o-monomial. Accordingly, the result is
not presented as a solution of Conjecture 2 in its unrestricted logical form.

\section{Computational checks}


The proofs above are symbolic. From the repository root, two exact Python
programs provide finite checks of the formulas:

\begin{itemize}
\item \texttt{loops/MATH-\allowbreak{}DIR-\allowbreak{}0048/verification/verify\_\allowbreak{}bar\_\allowbreak{}segre\_\allowbreak{}value\_\allowbreak{}polynomial.py}
reconstructs the overline-Segre
value-set polynomial for \(m=5,7,9,11\) and checks
\((c_1,c_2,c_3)=(0,0,1)\);
\item \texttt{loops/MATH-\allowbreak{}DIR-\allowbreak{}0049/verification/verify\_\allowbreak{}glynn\_\allowbreak{}cycle\_\allowbreak{}cover\_\allowbreak{}arithmetic.py}
checks the zero-weight formulas,
Kummer valuations, coefficient gap, and the first finite value-set
polynomials.

\end{itemize}

On Windows with the project environment, run:

\begin{quote}\footnotesize\ttfamily
.\textbackslash{}\allowbreak{}.venv\textbackslash{}\allowbreak{}Scripts\textbackslash{}\allowbreak{}python.exe\ \allowbreak{}loops\textbackslash{}\allowbreak{}MATH-\allowbreak{}DIR-\allowbreak{}0048\textbackslash{}\allowbreak{}verification\textbackslash{}\allowbreak{}verify\_\allowbreak{}bar\_\allowbreak{}segre\_\allowbreak{}value\_\allowbreak{}polynomial.py\par
\end{quote}

\begin{quote}\footnotesize\ttfamily
.\textbackslash{}\allowbreak{}.venv\textbackslash{}\allowbreak{}Scripts\textbackslash{}\allowbreak{}python.exe\ \allowbreak{}loops\textbackslash{}\allowbreak{}MATH-\allowbreak{}DIR-\allowbreak{}0049\textbackslash{}\allowbreak{}verification\textbackslash{}\allowbreak{}verify\_\allowbreak{}glynn\_\allowbreak{}cycle\_\allowbreak{}cover\_\allowbreak{}arithmetic.py\par
\end{quote}


These calculations test the implementation and small admissible parameters;
they are not substitutes for the general arguments.

\section*{AI-assisted research and writing disclosure}


AI tools were used to assist with exploratory computation, proof checking,
literature search, and drafting. Carptopus is the sole author and takes
responsibility for the mathematical claims and the final text.

\begin{thebibliography}{9}
\footnotesize
\raggedright

\bibitem{ding-tang}
C.~Ding and C.~Tang,
\emph{Infinite families of 3-designs from o-polynomials},
Advances in Mathematics of Communications \textbf{15} (2021), 557--573.
\href{https://doi.org/10.3934/amc.2020082}{doi:10.3934/amc.2020082}.

\bibitem{maschietti}
A.~Maschietti,
\emph{Difference set and hyperovals},
Designs, Codes and Cryptography \textbf{14} (1998), 89--98.
\href{https://doi.org/10.1023/A:1008264606494}{doi:10.1023/A:1008264606494}.

\bibitem{kolsch-kyureghyan}
L.~K\"olsch and G.~Kyureghyan,
\emph{The classifications of o-monomials and of 2-to-1 binomials are equivalent},
Designs, Codes and Cryptography \textbf{93} (2025), 961--970.
\href{https://doi.org/10.1007/s10623-024-01463-1}{doi:10.1007/s10623-024-01463-1}.

\bibitem{glynn}
D.~G.~Glynn,
\emph{Two new sequences of ovals in finite Desarguesian planes of even order},
in: Combinatorial Mathematics X, Lecture Notes in Mathematics \textbf{1036},
Springer, 1983, 217--229.
\href{https://doi.org/10.1007/BFb0071521}{doi:10.1007/BFb0071521}.

\end{thebibliography}

\end{document}
